\Needspace{7\baselineskip}
\Article{316}
Nach diesen Grundsätzen haben wir für alle Nenner der Form \(p^\mu\) unter \(1000\) eine Tafel der nötigen Perioden aufgestellt. Wir werden sie bei Gelegenheit vollständig oder auch noch weiter fortgesetzt veröffentlichen. Hier möge Tafel III, die nur bis \(100\) reicht, als Beispiel genügen; sie braucht kaum eine Erklärung. Für die Nenner, nach denen \(10\) eine Primitivwurzel ist, enthält sie die Perioden der Brüche mit dem Zähler \(1\), nämlich für \(7,17,19,23,29,47,59,61,97\). Für die übrigen gibt sie \(f\) Perioden zu den Zählern \(1,r,rr,\ldots r^{f-1}\), die durch die beigefügten Zahlen \((0),(1),(2)\) usw. unterschieden sind. Als Basis \(r\) wurde stets dieselbe Primitivwurzel wie in Tafel I gewählt. Daraus lässt sich also die Periode jedes Bruchs, dessen Nenner in dieser Tafel vorkommt, nach den Regeln des vorigen Artikels ermitteln, nachdem der Index des Zählers aus Tafel I berechnet wurde. Bei so kleinen Nennern können wir die Aufgabe aber ebenso leicht ohne Tafel I lösen. Durch gewöhnliche Division berechnen wir so viele Anfangsziffern der gesuchten Mantisse, wie nach Art.\ \(313\) nötig sind, um sie von allen anderen mit demselben Nenner zu unterscheiden; für Tafel III nicht mehr als \(2\). Dann durchsuchen wir sämtliche Perioden zum gegebenen Nenner, bis wir auf jene Anfangsziffern stoßen, die den Beginn der Periode zweifellos anzeigen. Wir müssen jedoch bemerken, dass diese Ziffern auch getrennt sein können: Die erste oder mehrere können das Ende einer Periode bilden, die übrige oder die übrigen deren Anfang.

\textit{Beispiel.} Gesucht ist die Periode von \(\frac{12}{19}\). Hier hat man nach dem Modul \(19\) aus Tafel I
\SourceLeaf{391}
die Beziehung \(\operatorname{ind.}12=2\operatorname{ind.}2+\operatorname{ind.}3=39\equiv3\; (\mathrm{mod}.\ 18)\), nach Art.\ \(57\). Da es in diesem Fall nur eine Periode zum Zähler \(1\) gibt, müssen ihre ersten drei Ziffern an das Ende gestellt werden. So entsteht die gesuchte Periode \(631578947368421052\). — Der Beginn der Periode hätte sich ebenso leicht aus den ersten beiden Ziffern \(63\) finden lassen.

Sucht man die Periode von \(\frac{45}{53}\), so gilt nach dem Modul \(53\) die Beziehung \(\operatorname{ind.}45=2\operatorname{ind.}3 +\operatorname{ind.}5=49\). Die Anzahl der Perioden ist hier \(4=f\), und es gilt \(49=12f+1\). Daher müssen von der mit \((1)\) bezeichneten Periode die ersten \(12\) Ziffern hinter die letzte gestellt werden. Die gesuchte Periode ist somit \(8490566037735\). Die Anfangsziffern \(84\) sind in diesem Fall in der Tafel getrennt.

Wir bemerken noch, dass sich mit Tafel III auch die Zahl finden lässt, die nach einem gegebenen Modul, der darin unter der Überschrift des Nenners vorkommt, einem gegebenen Index entspricht, wie wir in Art.\ \(59\) versprochen haben. Denn nach dem Vorangegangenen lässt sich offenbar die Periode des Bruchs finden, dessen Zähler, obwohl unbekannt, dem gegebenen Index entspricht. Es genügt, daraus so viele Anfangsziffern zu entnehmen, wie der Nenner Ziffern hat. Aus ihnen ergibt sich nach Art.\ \(313\) der Zähler, also die gesuchte Zahl zum gegebenen Index.

\Needspace{7\baselineskip}
\Article{317}
Nach dem Vorangegangenen lässt sich die Mantisse jedes Bruchs, dessen Nenner eine Primzahl oder eine Primzahlpotenz innerhalb der Grenzen der Tafel ist, ohne Rechnung auf beliebig viele Ziffern bestimmen. Mithilfe der Untersuchungen am Anfang dieses Abschnitts reicht der Umfang der Tafel aber viel weiter: Er umfasst alle Brüche, deren Nenner Produkte von Primzahlen oder Primzahlpotenzen innerhalb ihrer Grenze sind. Denn ein solcher Bruch lässt sich in andere zerlegen, deren Nenner diese Faktoren sind. Diese können auf beliebig viele Stellen in Dezimalbrüche umgewandelt werden; dann bleibt nur, sie zu addieren. Wir brauchen kaum zu bemerken, dass die letzte Ziffer der so entstehenden Summe kleiner als die richtige sein kann. Der Fehlbetrag kann aber offenbar nicht so viele Einheiten erreichen, wie einzelne Brüche addiert werden. Daher ist es zweckmäßig, diese auf einige Ziffern mehr zu berechnen, als man für den gegebenen Bruch an richtigen Ziffern benötigt. Als Beispiel betrachten wir den Bruch \(\frac{6099380351}{1271808720}=F\)\HistoricalNote{*)}{Dieser Bruch ist einer von denen, die der Quadratwurzel aus \(23\) so nahe wie möglich kommen; sein Überschuss ist kleiner als sieben Einheiten an der zwanzigsten Dezimalstelle.}, dessen Nenner das Produkt der Zahlen
\SourceLeaf{392}
\(16,\allowbreak 9,\allowbreak 5,\allowbreak 49,\allowbreak 13,\allowbreak 47,\allowbreak 59\) ist. Nach den obigen Regeln findet man \(F=1+\frac{11}{16}+\frac49+\frac45+\frac{22}{49}+\frac5{13}+\frac7{47}+\frac{52}{59}\). Diese einzelnen Brüche werden wie folgt in Dezimalbrüche umgewandelt:
\[\renewcommand{\arraystretch}{1.25}\begin{array}{r@{\;}c@{\;}l@{\quad}l@{\quad}l}
1&{}={}&1&&\\
\frac{11}{16}&{}={}&0,6875&&\\
\frac45&{}={}&0,8&&\\
\frac49&{}={}&0,4444444444&4444444444&44\\
\frac{22}{49}&{}={}&0,4489795918&3673469387&75\\
\frac5{13}&{}={}&0,3846153846&1538461538&46\\
\frac7{47}&{}={}&0,1489361702&1276595744&68\\
\frac{52}{59}&{}={}&0,8813559322&0338983050&84\\\hline
F&{}={}&4,7958315233&1271954166&17
\end{array}\]
Der Fehlbetrag dieser Summe gegenüber dem richtigen Wert ist sicher kleiner als fünf Einheiten an der letzten, zweiundzwanzigsten Stelle. Daher können die ersten einundzwanzig Ziffern dadurch nicht geändert werden. Setzt man die Rechnung auf mehr Ziffern fort, so erhält man anstelle der letzten beiden Ziffern \(17\) die Ziffern \(1893936\ldots\) — Auch ohne unseren Hinweis wird jeder sehen, dass dieses Verfahren, gewöhnliche Brüche in Dezimalbrüche umzuwandeln, besonders für den Fall geeignet ist, dass viele Dezimalziffern verlangt werden. Genügen wenige, so können gewöhnliche Division oder Logarithmen meist ebenso bequem angewandt werden.

\Needspace{28\baselineskip}
\Article{318}
Da nun die Umwandlung der Brüche, deren Nenner aus mehreren verschiedenen Primzahlen zusammengesetzt sind, auf den Fall zurückgeführt ist, dass der Nenner eine Primzahl oder eine Primzahlpotenz ist, fügen wir über ihre Mantissen nur wenige Bemerkungen hinzu. Enthält der Nenner weder den Faktor \(2\) noch \(5\), so besteht auch hier die Mantisse aus Perioden. Denn auch in diesem Fall erreicht man in der Reihe \(10,\allowbreak 100,\allowbreak 1000\) schließlich ein Glied, das nach dem Nenner der Einheit kongruent ist. Zugleich gibt dessen Exponent, der sich nach Art.\ \(92\) leicht bestimmen lässt, die vom Zähler unabhängige Länge der Periode an, sofern dieser zum Nenner teilerfremd ist. — Hat der Nenner dagegen die Form \(2^\alpha5^\beta N\), wobei \(N\) eine zu \(10\) teilerfremde Zahl ist und \(\alpha\) und \(\beta\) Zahlen sind, von denen wenigstens eine nicht \(0\) ist, so beginnt die Mantisse nach den ersten \(\alpha\) oder \(\beta\) Ziffern aus Perioden zu bestehen, je nachdem \(\alpha\) oder \(\beta\) größer ist. Diese stimmen in ihrer Länge mit den Perioden von Brüchen mit dem Nenner \(N\) überein.
\SourceLeaf{393}
Dies folgt sehr leicht daraus, dass sich jener Bruch in zwei andere mit den Nennern \(2^\alpha5^\beta\) und \(N\) zerlegen lässt, von denen der erste nach den ersten \(\alpha\) oder \(\beta\) Ziffern abbricht. — Wir könnten über diesen Gegenstand viele weitere Bemerkungen hinzufügen, besonders über Kunstgriffe, eine Tafel wie III möglichst schnell aufzustellen. Der Kürze wegen lassen wir sie hier umso lieber weg, als schon sowohl der angesehene Robertson am angeführten Ort als auch der angesehene Bernoulli (\textit{Nouv. Mém. de l'Ac. de Berlin} \(1771\ p.\ 273\)) vieles hierzu mitgeteilt haben.
